% Propietario conservado: /Users/ruben/Documents/New project/output/EXTREMA_MEDIA_RAZON_RECIPROCIDAD_ES_EN_20260918/ARTICULO/ES/source/nuclear/03.tex
\chapter{Recuperación del estado y de su dinámica}
\label{x_reciprocidad:nuclear:3}
\section{Representación recuperable del transporte}

Sea \(C\) la realización unitaria del transporte de memoria sobre \(\mathcal H\). La carta de nueve fases, su inclusión uniforme y su compresión son

\[
S_C(v_0,\ldots,v_8)=(Cv_8,v_0,\ldots,v_7),\qquad
Jv=(v,\ldots,v)/3,\qquad T=(8I+C)/9.
\]

Sea \(\eta=(I-JJ^*)S_CJ\). Su acción es

\[
\eta v=\frac1{27}\bigl(8(C-I)v,-(C-I)v,\ldots,-(C-I)v\bigr).
\]

La normalización de nueve copias determina \(1/3\); el reparto de la costura determina \(8/81\) en

\begin{equation}
I-T^*T=\eta^*\eta=\frac8{81}(I-C)^*(I-C).
\label{x_reciprocidad:nuc03:eq:1}
\end{equation}

Escribimos \(P=\operatorname{proj}\ker(I-C)\). El teorema precedente prueba la convergencia fuerte \(T^N\to P\). Por tanto,

\begin{equation}
\mathcal Vv=(Pv,\eta v,\eta Tv,\eta T^2v,\ldots)
\quad\text{en}\quad
\mathcal R=P\mathcal H\oplus\ell^2(\mathbb N_0;\mathcal H^9)
\label{x_reciprocidad:nuc03:eq:2}
\end{equation}

es una isometría. Denotamos su imagen cerrada por \(\mathcal M\).

En el espacio de registros definimos el desplazamiento de bloques

\begin{equation}
\mathcal B(a,m_0,m_1,m_2,\ldots)=(a,m_1,m_2,\ldots).
\label{x_reciprocidad:nuc03:eq:3}
\end{equation}

El primer componente conserva el sector fijo; los demás componentes son los registros vectoriales completos, incluidas sus correlaciones.

\begin{theorem}
\label{x_reciprocidad:nuc03:transporte}
El transporte original tiene la representación exacta

\begin{equation}
\boxed{\mathcal V C=(9\mathcal B-8I)\mathcal V.}
\label{x_reciprocidad:nuc03:eq:4}
\end{equation}

La restricción \(\mathcal C=(9\mathcal B-8I)|_{\mathcal M}\) es unitaria sobre \(\mathcal M\), y

\begin{equation}
C=\mathcal V^*\mathcal C\mathcal V.
\label{x_reciprocidad:nuc03:eq:5}
\end{equation}
\end{theorem}

\begin{proof}
Como \(PT=P\), la definición~\eqref{x_reciprocidad:nuc03:eq:2} da componente a componente

\[
\mathcal VT v=(Pv,\eta Tv,\eta T^2v,\ldots)=\mathcal B\mathcal Vv.
\]

La identidad algebraica \(C=9T-8I\) implica \eqref{x_reciprocidad:nuc03:eq:4}. La isometría \(\mathcal V:\mathcal H\to\mathcal M\) es sobreyectiva por definición de su imagen. La igualdad \eqref{x_reciprocidad:nuc03:eq:4} identifica la restricción con \(\mathcal V C\mathcal V^{-1}\), que es unitaria porque \(C\) lo es. Esto demuestra además que \(\mathcal C\mathcal M=\mathcal M\) y \eqref{x_reciprocidad:nuc03:eq:5}.
\end{proof}

La restricción al subespacio de registros compatibles es esencial: \(9\mathcal B-8I\) sobre todo \(\mathcal R\) no es unitario. El teorema caracteriza el transporte sobre registros que proceden efectivamente de una misma trayectoria lineal de la carta documentada.

\begin{corollary}
Para \(v,w\in\mathcal H\) y \(k\in\mathbb Z\),

\begin{equation}
\langle C^kv,w\rangle
=\langle\mathcal C^k\mathcal Vv,\mathcal Vw\rangle.
\label{x_reciprocidad:nuc03:eq:6}
\end{equation}

Así se conservan todos los momentos cruzados de la dinámica realizada. Las potencias negativas pertenecen a la inversa de \(\mathcal C\) en \(\mathcal M\), no a una inversión del desplazamiento unilateral sobre el espacio ambiente.
\end{corollary}

\begin{corollary}
Para todo observable acotado \(A\), su realización en la imagen es \(\widehat A=\mathcal V A\mathcal V^*|_{\mathcal M}\). Se preservan productos, adjuntos y conmutadores. En particular,

\begin{equation}
[A,C]=0\quad\Longleftrightarrow\quad[\widehat A,\mathcal C]=0.
\label{x_reciprocidad:nuc03:eq:7}
\end{equation}
\end{corollary}

La prueba consiste en usar \(\mathcal V^*\mathcal V=I\) y que \(\mathcal V\mathcal V^*\) es la identidad de \(\mathcal M\). La conservación de simetría queda expresada mediante el transporte íntegro de sus operadores, además de mediante una igualdad de normas.

\section{Reconstrucción explícita y profundidad del registro}

Sea \(m_j=\eta T^jv\). La identidad telescópica da para cada \(N\)

\begin{equation}
\boxed{v=T^{*N}T^Nv+\sum_{j=0}^{N-1}T^{*j}\eta^*m_j.}
\label{x_reciprocidad:nuc03:eq:8}
\end{equation}

Al pasar al límite fuerte,

\begin{equation}
\boxed{v=Pv+\sum_{j=0}^{\infty}T^{*j}\eta^*m_j.}
\label{x_reciprocidad:nuc03:eq:9}
\end{equation}

La reconstrucción parcial con terminal fijo,
\(v_N=Pv+\sum_{j<N}T^{*j}\eta^*m_j\), tiene error exacto

\begin{equation}
v-v_N=T^{*N}T^N(I-P)v.
\label{x_reciprocidad:nuc03:eq:10}
\end{equation}

La serie converge para cada estado. Si el registro completo presenta un error \(\delta r\) en norma de \(\mathcal R\), su reconstrucción mediante \(\mathcal V^*\) cambia en una norma a lo sumo \(\|\delta r\|\), porque \(\|\mathcal V^*\|=1\). Ésta es una estabilidad uniforme de la lectura completa.

\subsection{Diferencia entre profundidad finita e ilimitada}

Para el desplazamiento bilateral \(Ce_k=e_{k+1}\) en \(\ell^2(\mathbb Z)\), se tiene \(P=0\). Considérese

\[
v_L=L^{-1/2}\sum_{k=1}^{L}e_k.
\]

Entonces \(\|(I-C)v_L\|^2=2/L\). La conmutación de \(T\) con \(I-C\), la contracción de \(T\) y \eqref{x_reciprocidad:nuc03:eq:1} implican

\begin{equation}
\sum_{j=0}^{N-1}\|\eta T^jv_L\|^2
\leq\frac{16N}{81L}.
\label{x_reciprocidad:nuc03:eq:11}
\end{equation}

Para cualquier profundidad fija \(N\), el lado derecho tiende a cero cuando \(L\to\infty\), aunque \(\|v_L\|=1\). El registro de memoria de profundidad fija, separado del terminal, carece de una cota inferior uniforme. En cambio, el registro ilimitado tiene norma exactamente uno para cada \(v_L\).

La diferencia corresponde a dos cuantificadores distintos: cada estado queda conservado al completar su registro; una profundidad prefijada no recupera uniformemente todos los estados sin su terminal. La conservación a profundidad arbitraria posee aquí una expresión funcional exacta.

\subsection{Información vectorial e intensidad}

El teorema conserva \(m_j\), no únicamente \(\|m_j\|^2\). Por ejemplo, en la realización bilateral, \(v=e_0\) y \(w=e_1\) tienen idénticas intensidades en cada profundidad porque el desplazamiento conmuta con \(T\) y transporta \(\eta\) unitariamente. Sus registros vectoriales son diferentes y \eqref{x_reciprocidad:nuc03:eq:9} recupera los dos estados distintos. Las intensidades sirven para el balance; los registros con sus correlaciones sirven además para la reconstrucción.

\section{Naturalidad de la recuperación bajo refinamiento}

Sean \(\iota_n:\mathcal H_n\to\mathcal H_{n+1}\) las isometrías de refinamiento obtenidas de las emisiones comunes, y \(C_n\) realizaciones unitarias con

\begin{equation}
C_{n+1}\iota_n=\iota_nC_n.
\label{x_reciprocidad:nuc03:eq:12}
\end{equation}

La relación vale también para los adjuntos: multiplicar por \(C_{n+1}^*\) y \(C_n^*\) da \(C_{n+1}^*\iota_n=\iota_nC_n^*\). Por tanto las imágenes son reductoras. Las fórmulas polinómicas de \(T_n\), \(\eta_n\) y el límite fuerte de \(T_n^N\) dan

\begin{equation}
T_{n+1}\iota_n=\iota_nT_n,\quad
\eta_{n+1}\iota_n=\iota_n^{\oplus9}\eta_n,\quad
P_{n+1}\iota_n=\iota_nP_n.
\label{x_reciprocidad:nuc03:eq:13}
\end{equation}

Sea \(\mathfrak I_n\) el refinamiento que aplica \(\iota_n\) al terminal y \(\iota_n^{\oplus9}\) a cada registro. Entonces

\begin{equation}
\boxed{\mathcal V_{n+1}\iota_n=\mathfrak I_n\mathcal V_n,\qquad
\mathcal B_{n+1}\mathfrak I_n=\mathfrak I_n\mathcal B_n.}
\label{x_reciprocidad:nuc03:eq:14}
\end{equation}

\begin{proof}
Cada componente de la primera igualdad es una identidad de \eqref{x_reciprocidad:nuc03:eq:13}, iterada a la potencia correspondiente de \(T_n\). La segunda igualdad afirma que refinar cada bloque conmuta con desplazar el índice de bloques, lo cual se comprueba directamente. Ambas aplicaciones son acotadas, de modo que las igualdades pasan de los registros de soporte finito a su completación.
\end{proof}

Por \eqref{x_reciprocidad:nuc03:eq:4} y \eqref{x_reciprocidad:nuc03:eq:14}, la reconstrucción de la dinámica también conmuta con el refinamiento. Se mantienen distintos el índice de profundidad genealógica \(n\) y el de compresiones sucesivas \(j\). La igualdad expresa su compatibilidad cuando actúan mediante los cuadrados documentados, en lugar de identificarlos nominalmente.

\section{Dos regímenes exactos de distribución de memoria}

Para un estado unitario \(v\) con \(Pv=0\), definimos

\begin{equation}
a_N=\|T^Nv\|^2,\qquad b_N=\|\eta T^Nv\|^2=a_N-a_{N+1}.
\label{x_reciprocidad:nuc03:eq:15}
\end{equation}

Se tiene \(a_0=1\), \(a_N\to0\), \(b_N\geq0\) y \(\sum b_N=1\). Los pesos \(b_N\) describen la distribución entre profundidades del registro de memoria. Este índice cuenta compresiones; su identificación con una duración física requeriría aplicar el lector temporal correspondiente.

\subsection{Transporte excepcional de orden tres}

El operador ortogonal fijo-libre \(g_c\) satisface \(g_c^2+g_c+I=0\). Para la composición explícita \(C=g_c\), el desarrollo precedente obtuvo

\begin{equation}
T^*T=\frac{19}{27}I,
\qquad a_N=\left(\frac{19}{27}\right)^N,
\qquad b_N=\frac8{27}\left(\frac{19}{27}\right)^N.
\label{x_reciprocidad:nuc03:eq:16}
\end{equation}

La profundidad media, contando el primer registro como profundidad uno, es

\begin{equation}
\sum_{N\geq0}(N+1)b_N=\frac{27}{8}.
\label{x_reciprocidad:nuc03:eq:17}
\end{equation}

Es un resultado de esa composición de la carta nonádica con el transporte excepcional ya construido; conserva el dominio y no identifica por su orden a \(g_c\) con la holonomía temporal completa.

\subsection{Registro bilateral de vueltas}

Para \(C e_k=e_{k+1}\) y \(v=e_0\), el binomio determina todos los términos:

\begin{equation}
T^Ne_0=9^{-N}\sum_{k=0}^{N}\binom{N}{k}8^{N-k}e_k,
\quad
a_N=81^{-N}\sum_{k=0}^{N}\binom{N}{k}^{\!2}64^{N-k}.
\label{x_reciprocidad:nuc03:eq:18}
\end{equation}

Esta fórmula es exacta a cualquier profundidad. La representación circular posterior proporciona

\begin{equation}
a_N=\frac1{2\pi}\int_{-\pi}^{\pi}
\left(\frac{65+16\cos\theta}{81}\right)^N\,d\theta.
\label{x_reciprocidad:nuc03:eq:19}
\end{equation}

\begin{theorem}
En esta realización bilateral,

\begin{equation}
\boxed{a_N\sim\frac9{4\sqrt{2\pi}}N^{-1/2},\qquad
b_N\sim\frac9{8\sqrt{2\pi}}N^{-3/2}.}
\label{x_reciprocidad:nuc03:eq:20}
\end{equation}
\end{theorem}

\begin{proof}
Escribimos \(r(\theta)=(65+16\cos\theta)/81\) y \(q=8/81\). En torno a cero,

\[
r(\theta)=1-q\theta^2+O(\theta^4),\qquad
\log r(\theta)=-q\theta^2+O(\theta^4).
\]

Fijado un entorno suficientemente pequeño, existen constantes positivas \(c_1,c_2\) que acotan \(-\log r(\theta)\) entre \(c_1\theta^2\) y \(c_2\theta^2\). Fuera de ese entorno, \(r\leq\rho<1\). El cambio \(x=\sqrt N\theta\), la cota gaussiana y la convergencia dominada dan

\[
\lim_{N\to\infty}\sqrt N\,a_N
=\frac1{2\pi}\int_{\mathbb R}e^{-qx^2}\,dx
=\frac1{2\sqrt{\pi q}}=\frac9{4\sqrt{2\pi}}.
\]

Para el segundo límite se aplica el mismo procedimiento directamente a

\[
b_N=\frac1{2\pi}\int_{-\pi}^{\pi}r(\theta)^N(1-r(\theta))\,d\theta.
\]

En la variable \(x\), \(N(1-r(x/\sqrt N))\to qx^2\), dominado localmente por una constante por \(x^2\); fuera del entorno la contribución es exponencialmente pequeña incluso tras multiplicar por \(N^{3/2}\). Por tanto

\[
\lim_{N\to\infty}N^{3/2}b_N
=\frac q{2\pi}\int_{\mathbb R}x^2e^{-qx^2}\,dx
=\frac1{4\sqrt{\pi q}}=\frac9{8\sqrt{2\pi}}.
\]

Esta prueba del segundo límite usa su integral propia, sin diferenciar una equivalencia asintótica no controlada.
\end{proof}

Por Tonelli y la identidad de colas de una sucesión positiva,

\begin{equation}
\sum_{N\geq0}(N+1)b_N=\sum_{N\geq0}a_N=+\infty.
\label{x_reciprocidad:nuc03:eq:21}
\end{equation}

Toda la norma se conserva en memoria, pero su profundidad media diverge en este estado bilateral. La ley de reparto depende de la representación efectiva del transporte: el sector excepcional de orden tres posee tasa geométrica, mientras el registro bilateral aquí considerado posee tasa algebraica.

\section{Composición con la identidad de Euler y la incidencia}

La identidad de Euler extendida de la sección~\ref{x_reciprocidad:euler-trit-generadores} realiza los tres regímenes mediante el Coxeter de \(A_2\), un transporte nilpotente y \(F_{\rm av}^2\), con sus antecedentes declarados. Sus generadores normalizados satisfacen \(J_\tau^2=-\tau I\) y tienen traza cero. En sus cartas bidimensionales,

\[
C_\tau(t)=\exp(tJ_\tau),\qquad
\Omega=\begin{pmatrix}0&1\\-1&0\end{pmatrix},\qquad
C_\tau(t)^\mathsf T\Omega C_\tau(t)=\Omega.
\]

La traza cero da \(\det C_\tau(t)=1\), y la identidad \(M^\mathsf T\Omega M=(\det M)\Omega\) en dimensión dos prueba la última igualdad. Para

\[
T_\tau=\frac{8I+C_\tau}{9},\qquad
D_\tau=\frac{\sqrt8}{9}(C_\tau-I),
\]

se obtiene una ley de conservación orientada común:

\begin{equation}
\boxed{\Omega=T_\tau^\mathsf T\Omega T_\tau+D_\tau^\mathsf T\Omega D_\tau.}
\label{x_reciprocidad:nuc03:eq:22}
\end{equation}

\begin{proof}
Al expandir los dos términos, las contribuciones cruzadas de \(C_\tau\) tienen coeficientes opuestos. Quedan \(72\Omega/81+9C_\tau^\mathsf T\Omega C_\tau/81=\Omega\).
\end{proof}

Esta forma alternante corresponde al área orientada del plano canónico. Su conservación permite comparar los tres regímenes respetando sus diferencias. Para los representantes concretos \(C\) de Coxeter, \(I+N\) y \(F_{\rm av}^2\), respectivamente,

\begin{equation}
(\det T,\det D)=
\left(\frac{19}{27},\frac8{27}\right),\qquad
(1,0),\qquad
\left(\frac{89}{81},-\frac8{81}\right).
\label{x_reciprocidad:nuc03:eq:23}
\end{equation}

La suma es uno en los tres casos. En el parabólico, \(D\) es de rango uno y puede conservar información aunque su área sea cero. En el hiperbólico, el complemento tiene orientación opuesta; el balance es de área firmada. El teorema~\ref{x_reciprocidad:nuc03:transporte} utiliza una realización unitaria y positiva: se conservan expresamente las hipótesis de cada resultado.

El desarrollo de acción y forma variable prueba la versión entre cartas diferentes, su integración sobre curvas cerradas y su compatibilidad con refinamiento. El desarrollo de incidencia compone la reconstrucción con \(F(k)=(\mu(k),\mathcal I P_0k/6)\), manteniendo las doce coordenadas de \(K\); transporta también los operadores y sus marcos locales. Así, una misma identidad de conservación admite realizaciones tipadas de norma, incidencia y acción orientada.
